\section{Fundamentação rastreável e limites de transferência}

Esta síntese exploratória relaciona planejamento, interfaces, funcionamento local, engenharia e integridade da escrita. Foram consultados principalmente resumos e metadados; não foi realizada avaliação sistemática do risco de viés de cada estudo. As fontes sustentam os enunciados delimitados abaixo, sem validar o EstudaOffline nem comprovar aprendizagem. O registro de busca e seleção integra o protocolo do projeto \citep{projeto}.

\subsection{Planejamento e autorregulação da aprendizagem}

Revisão identificou uso continuado de instrumentos de sala presencial para medir autorregulação em ambientes digitais e lacunas em análises de registros. Período 2008–2018; medir atividades digitais não equivale a medir aprendizagem. \citep{education_araka2020}.


Revisão teórica discute suporte móvel oportuno a planejamento, monitoramento, reflexão e estratégias de aprendizagem. Não é ensaio de um planejador; oportunidades conceituais não são eficácia demonstrada. \citep{education_baars2022}.


Levantamento com 142 universitários encontrou associação negativa entre autorregulação e procrastinação acadêmica. Questionários e regressão não demonstram causalidade nem validade local para alunos do ADS. \citep{education_brahma2023}.


Em 73 universitários, condição combinando treinamento online e diários móveis apresentou maiores melhorias em componentes autorregulatórios. Intervenção combinada impede atribuir o resultado ao diário isolado ou a qualquer aplicativo. \citep{education_broadbent2020}.


Em duas turmas de 22 participantes, abordagem móvel gamificada melhorou desempenho, definição de metas e reflexão, mas não monitoramento. Amostra pequena; gamificação e instrução diferem de organizador local de tarefas. \citep{education_chen2024}.


Em 8.907 universitários durante a pandemia, motivação e procrastinação relacionaram-se à gestão do tempo e regulação do esforço. Contexto de emergência e relações observacionais não generalizam automaticamente para ensino técnico. \citep{education_forstervold2022}.


Em 1.216 universitários, autoeficácia autorregulatória e procrastinação apresentaram relações recíprocas; gestão do tempo foi dificuldade recorrente. Cursos assíncronos durante COVID-19; não testa aplicativo offline. \citep{education_han2023}.


Planejamento por smartphone e agenda associou-se a habilidades autorregulatórias; modelos apresentaram associações diferentes com notas. Associação negativa de smartphone e positiva de agenda sob controle estatístico não é causalidade nem avalia este aplicativo. \citep{education_hartley2022}.


Duas pesquisas-ação descrevem regulação não linear e construção compartilhada de conhecimento em aprendizagem móvel colaborativa. Contexto colaborativo qualitativo não valida planejador individual nem produz estimativa causal. \citep{education_lim2019}.


Em 450 universitários, estratégias metacognitivas e gestão do tempo foram examinadas como preditores de procrastinação. Desenho transversal não estabelece causalidade; contexto universitário específico. \citep{education_limone2020}.


Revisão de 38 estudos descreveu relação complexa entre aprendizagem móvel e autorregulação e recomendou integração curricular intencional. Síntese heterogênea de 2007–2019 não demonstra efeito de todo app móvel. \citep{education_palalas2020}.


Revisão de 107 estudos identificou planejamento, metas, priorização e organização de tarefas entre estratégias relevantes. Inclui educação e trabalho, artigos e dissertações; não demonstra efeito do MVP nem padroniza medida universal. \citep{education_patzak2025}.


Síntese de 31 revisões e meta-análises recomenda suportes autorregulatórios integrados ao longo do ciclo de aprendizagem. Revisões com contextos diversos; recomendação de desenho não constitui validação deste MVP. \citep{education_prasse2024}.


Em 272 universitários iranianos de inglês, maior autorregulação e controle atencional associaram-se a menor procrastinação. Amostra por conveniência e autorrelato; registro Consensus omitiu terceiro autor, corrigido pelo editor; 2024 online e 2026 fascículo. \citep{education_taghavi2024}.


Em 239 universitários de inglês como língua estrangeira, gestão do tempo e regulação do esforço associaram-se negativamente à procrastinação. Modelo de associações não demonstra eficácia de um planejador; população específica. \citep{education_tao2025}.


Em levantamento com 446 universitários, gestão do tempo contribuiu à predição de procrastinação além de variáveis motivacionais e estratégicas. Autorrelato e regressão não estabelecem efeito causal de um aplicativo. \citep{education_wolters2017}.


Revisão situa a gestão do tempo nas fases de antecipação, execução e avaliação da aprendizagem autorregulada. Síntese conceitual não comprova eficácia de aplicativo específico. \citep{education_wolters2021}.


Estudo identificou perfis ativo, desengajado e utilitário em app de prática de recuperação; autoeficácia e gestão do tempo predisseram notas. Perfis e associações não são efeito causal de uso de aplicativo; app inclui prática de recuperação ausente neste MVP. \citep{education_wong2026}.


Meta-análise encontrou efeito positivo moderado de intervenções de autorregulação no desempenho em ambientes online e híbridos. Intervenções e níveis escolares heterogêneos; efeito agrupado não valida EstudaOffline; 2022 online e 2023 fascículo. \citep{education_xu2022}.


Experimento com 60 graduandos encontrou melhores resultados e autorregulação em leitura de inglês apoiada por mecanismo móvel autorregulatório. Intervenção específica de leitura; 2016 online e 2018 fascículo; não transfere efeito ao EstudaOffline. \citep{education_zheng2016}.


\subsection{Acessibilidade, usabilidade e desigualdade digital}

Cobertura de rede, dispositivo e condições socioeconômicas afetam participação em ensino síncrono. Netnografia de tweets, com seleção de usuários de rede social; não representa todos os estudantes. \citep{azionya2021}.


Avaliação por pessoas cegas e surdas pode revelar diferenças na facilidade de uso de uma plataforma adaptada. Protótipo e população específicos; percepções não demonstram eficácia pedagógica ou acessibilidade universal. \citep{batanero2021}.


A revisão examina a formação docente para produção de materiais e ambientes digitais inclusivos; adotar apoio institucional no EstudaOffline é uma proposta de projeto. Revisão de 16 estudos; predominam entrevistas e surveys, com poucos indicadores objetivos. Não avalia EstudaOffline. \citep{bong2021}.


Acessibilidade, usabilidade e desenho universal para aprendizagem devem ser considerados em conjunto. Artigo conceitual com exemplos; não estima efeito causal nem certifica acessibilidade de um aplicativo. \citep{choi2024}.


Há lacunas na avaliação de efetividade de e-learning acessível a pessoas com dificuldades cognitivas. Revisão de estudos heterogêneos; resumo do editor consultado, texto integral não inspecionado. \citep{cinquin2019}.


A revisão organiza critérios de interface e sua relação com carga cognitiva; considerar tamanho de tela e desenho da interface no EstudaOffline é uma inferência de projeto. Revisão limitada a 2020–2022; não prova redução de carga cognitiva por uma interface específica. \citep{cob2023}.


O estudo avaliou recursos de voz e língua de sinais com estudantes e docentes; propor avaliação participativa do EstudaOffline é uma inferência de projeto. Estudo de métodos mistos de satisfação; não transfere resultados ao EstudaOffline nem à educação técnica brasileira. \citep{farhan2020}.


A avaliação de acessibilidade educacional deve considerar barreiras atitudinais além de aspectos técnicos. A síntese não permite afirmar que melhorias técnicas removem discriminação ou todas as barreiras. \citep{gobbo2023}.


Qualidade da conexão e condições de acesso integram a divisão digital no ensino online. Survey de estudantes chineses no contexto pandêmico; não autoriza inferência causal nem generalização direta ao Piauí. \citep{guo2022}.


Acesso à internet e habilidades digitais podem limitar continuidade educacional em universidades. Entrevistas por email podem excluir as pessoas com pior acesso; contexto paquistanês e pandêmico. \citep{jamil2023}.


Usabilidade de aplicações de aprendizagem móvel exige avaliação explícita e não pode ser presumida pela tecnologia. 23 publicações relevantes na revisão; publicado online em 2017, edição de 2018; não mede EstudaOffline. \citep{kumar2017}.


Diretrizes de usabilidade móvel devem ser atualizadas a partir de problemas observados. Síntese de 17 estudos e validação das diretrizes; condições e efeitos não podem ser transportados automaticamente. \citep{kumarguideline2019}.


Testes com usuários, observação e questionários são utilizados na avaliação de aprendizagem móvel. Mapeamento de 51 trabalhos de uma base; muitos testes controlados, limitando transferência para uso real. \citep{kumarmapping2024}.


Entrevistas com docentes e administradores descrevem estratégias de enfrentamento da divisão digital; combinar conectividade acessível, suporte e formação no projeto é uma proposta informada por esses relatos. 35 entrevistas em cinco universidades de Gana; relato de estratégias não é ensaio de eficácia causal. \citep{kumiyeboah2023}.


Existe pesquisa anterior sobre AVA com operação offline e uso em baixa conectividade, tornando inadequada uma alegação genérica de novidade de aplicativo offline educacional. Educa Offline é um AVA distinto de um planejador; o resumo e metadados não sustentam equivalência arquitetural, comparação de desempenho ou eficácia do EstudaOffline. SciSpace omitiu terceiro autor, corrigido com fonte primária. \citep{linhalis2026}.


Inclusão digital educacional depende de acesso, uso, fatores sociais e suporte organizacional. Revisão de 77 estudos em duas décadas; heterogeneidade e publicação online/volume posterior impedem uma estimativa única. \citep{martin2024}.


Acesso, práticas culturais e agência devem compor a análise de aprendizagem com conectividade limitada. Estudo em cinco países e no lockdown; gênero e responsabilidades domésticas dependem do contexto. \citep{mathrani2021}.


Facilidade de uso e acessibilidade integram fatores considerados na adoção de aprendizagem móvel. Revisão de 161 artigos, majoritariamente quantitativos; adoção e aprendizagem não são equivalentes. \citep{naveed2023}.


O artigo revisa padrões de interfaces móveis; aplicar sua discussão ao planejamento escolar do EstudaOffline é uma inferência de projeto. Revisão de interfaces móveis em geral; sua aplicação a planejamento escolar é uma inferência, não resultado direto. \citep{punchoojit2017}.


Prontidão digital dos estudantes e das escolas são dimensões distintas da desigualdade educacional. Dados e contextos internacionais; não demonstra que uma PWA resolve desigualdade estrutural. \citep{werfhorst2022}.


\subsection{Persistência local, cache, segurança e privacidade}

O estudo mostra que XSS na mesma origem pode manipular Cache API e ampliar o impacto sobre conteúdo armazenado. Cenários de ataque específicos; não demonstra vulnerabilidade no EstudaOffline nem estado atual de todos os navegadores. \citep{OS01}.


A sistematização reproduz ataques ligados a service workers e analisa medidas de mitigação dos navegadores. Panorama de 2022; mitigação atual exige consulta e teste específico. \citep{OS02}.


Os autores demonstram ataques de inferência de histórico explorando isolamento de service workers e cache. Resultados relativos às versões investigadas; não presumir que os ataques persistem em navegadores atuais. \citep{OS03}.


Experimentos mostram limitações dos controles e salvaguardas de cache nos navegadores examinados. Cache HTTP não equivale à Cache API; resultados históricos não certificam navegadores atuais. \citep{OS06}.


LoRe combina programação reativa e verificação para identificar interações que violam invariantes e coordená-las quando necessário. Garantias dependem do modelo, das propriedades especificadas e do compilador; não se transferem ao Flutter localStorage. \citep{OS07}.


Forking histories preserva restrições por histórico e expõe conflitos para reconciliação pelo usuário. Proposta de sincronização; não é resultado de adoção educacional nem regra necessária para app sem colaboração. \citep{OS08}.


O trabalho compara opções de criptografia para Web Storage e relata alternativas com impacto de desempenho limitado no cenário avaliado. Criptografia local não elimina XSS nem valida gestão de chaves; dados e implementação do MVP não foram avaliados. \citep{OS09}.


BeauForT propõe consenso e sincronização no navegador para tolerar réplicas bizantinas e redes instáveis. Sistema descentralizado com múltiplas partes; não constitui proteção pronta para aplicativo local sem sincronização. \citep{OS10}.


A análise de 137 apps infantis e entrevistas com 12 desenvolvedores identificaram barreiras ao uso de SDKs favoráveis à privacidade. Estudo qualitativo e ferramenta de apoio; não demonstra que remover SDKs garante segurança integral. \citep{OS13}.


A análise de 42 apps Android de cuidado infantil e seus backends identificou rastreadores e riscos de exposição de dados. Contexto de creches e backends móveis; não extrapolar vulnerabilidades para app local com dados sintéticos. \citep{OS14}.


Oficinas de codesign orientaram KOALA Hero para tornar riscos de tratamento de dados compreensíveis a crianças. Os autores deixam avaliação empírica futura em aberto; não evidencia mudança comprovada de comportamento. \citep{OS15}.


Análise estática e dinâmica de 20.195 apps identificou práticas de rastreamento e permissões que divergiam de políticas de proteção infantil. Amostra Google Play e políticas do período; não mede riscos de uma PWA específica nem todos os ecossistemas. \citep{OS17}.


A revisão selecionou 21 de 64 trabalhos e apontou lacunas na literatura sobre segurança de apps educacionais. Revisão secundária; apenas resumo consultado, sem auditoria de busca ou qualidade de cada estudo. \citep{OS18}.


\subsection{Engenharia, testes e reprodutibilidade}

A revisão de 103 estudos distingue testes funcionais e não funcionais e identifica desafios do ecossistema Android. Foco Android e literatura até 2016; não valida Flutter web nem este aplicativo. \citep{eng01}.


A revisão de 56 artigos identifica necessidades de adaptação dos frameworks de automação ao domínio móvel. Arquitetura MATF proposta, com implementação futura; não comprova eficácia do MVP. \citep{eng02}.


O mapeamento de 131 artigos organiza técnicas, atividades de teste e formas de avaliação, com predominância de Android. Estudo secundário; o ano online 2018 difere do volume de 2019. Não transferir resultados Android diretamente à web. \citep{eng03}.


O estudo terciário organiza 24 revisões sobre engenharia de aplicativos e identifica lacunas em manutenção e plataformas cruzadas. Síntese de revisões; não demonstra qualidade ou efeito pedagógico do aplicativo. \citep{eng04}.


O mapeamento de 79 estudos identifica necessidade de elicitar requisitos de teste cedo e avaliar em contextos reais. Mapeamento secundário de literatura anterior; não substitui ensaios atuais no navegador. \citep{eng05}.


O mapeamento de 56 estudos caracteriza técnicas e ferramentas de teste dinâmico de requisitos não funcionais. Classifica trabalhos existentes; segurança e confiabilidade do MVP exigem testes específicos. \citep{eng06}.


O artigo propõe explicitar o artefato e o objetivo de pesquisa para tornar mais precisa a investigação em design science. Validação metodológica em artigos DSR; não é medida de eficácia educacional. \citep{eng07}.


A reavaliação de um método para estudos de repositórios mostra utilidade de caracterizar a reprodutibilidade além da disponibilidade. Escopo de mineração de repositórios; não garante reprodução do MVP. \citep{eng08}.


A revisão e reexecução de modelos de deep learning mostram que compartilhar código não basta para reproduzir resultados. Escopo deep learning; o aplicativo não usa IA. Relevância limitada a documentação e pacotes reprodutíveis. \citep{eng10}.


O mapeamento de 40 estudos identifica dependência do domínio nas ferramentas e práticas de reprodução experimental. Ferramentas dependem do domínio; não autoriza generalização direta para estudantes. \citep{eng11}.


Sete replicações externas identificam dificuldades relacionadas a instruções pouco claras e defeitos nos artefatos. Experimento em sistemas configuráveis; aplicar como orientação de documentação, não como resultado do MVP. \citep{eng12}.


A revisão de replicações identifica relato incompleto e problemas para interpretar resultados confirmatórios. Escopo previsão de esforço e pair programming; relato do resumo não estabelece validade de qualquer replicação. \citep{eng13}.


A revisão de 69 artigos organiza ferramentas e desafios de integração, entrega e implantação contínuas. Revisão até junho de 2016; resultados exigem contextualização em pipelines atuais. \citep{eng14}.


A revisão de 101 estudos empíricos identifica benefícios e desafios técnicos associados à integração contínua. Estudos heterogêneos; benefícios de CI não equivalem a ausência de bugs ou eficácia pedagógica. \citep{eng15}.


O estudo de três empresas mostra que a implementação de CI varia com o contexto e pode criar restrições ao feedback. Três empresas; generalização limitada e ausência de avaliação deste pipeline. \citep{eng16}.


O FEDS separa avaliações formativas e somativas, artificiais e naturalistas, e orienta o planejamento dos episódios. Framework metodológico com exemplos; não certifica o artefato nem substitui avaliação com usuários. \citep{eng17}.


A proposta de estrutura de artefatos para experimentos humanos organiza elementos alinhados às políticas de badges ACM. Exige validação externa adicional; um repositório estruturado não recebe automaticamente badge ACM. \citep{eng18}.


O MEDS propõe planejar objetivos, estratégias, propriedades e episódios de avaliação para ajustar rigor e recursos em DSR. Método orientador; não demonstra qualidade do MVP nem efeitos nos estudantes. \citep{eng20}.


\subsection{Integridade, atribuição e revisão da escrita}

PRISMA 2020 fornece checklist de 27 itens e fluxogramas para relatar revisões sistemáticas de modo transparente. Diretriz de relato, não valida qualidade metodológica nem transforma esta busca exploratória em revisão sistemática. \citep{SU01}.


O artigo de explicação detalha a justificativa e exemplifica recomendações de cada item PRISMA 2020. Publicação complementar à mesma diretriz: fonte distinta, não estudo independente de eficácia. \citep{SU02}.


PRISMA-S apresenta 16 itens para descrever fontes e procedimentos de busca de forma reproduzível. Checklist desenvolvido por consenso; completude do relato não garante exaustividade ou ausência de viés. \citep{SU03}.


O manifesto propõe melhorar métodos, transparência, reprodução, avaliação e incentivos na pesquisa. Proposta baseada em evidências variadas; adoção ampla ainda requer avaliação iterativa e não certifica o MVP. \citep{SU04}.


A revisão de 239 estudos identifica avanços na detecção de plágio e lacunas de avaliação rigorosa dos sistemas. Literatura de 2013–2018; detecção computacional não substitui julgamento contextual nem garante originalidade. \citep{SU05}.


A avaliação de 636 referências geradas por GPT-3.5 e GPT-4 encontrou referências fabricadas e erros em referências reais. Modelos e prompts de abril de 2023; taxas não são estimativas para modelos atuais nem para esta pesquisa com verificação externa. \citep{SU06}.


Nos 30 textos médicos gerados por GPT-3.5, os autores encontraram referências fabricadas ou com metadados incorretos. 115 referências e prompts médicos específicos; não generalizar taxas para todos os gêneros ou modelos. \citep{SU07}.


AcPgChecker compara documentos offline com similaridade de cosseno e um limiar para sinalizar possível plágio. Resumo não apresenta métricas completas; limiar de similaridade não constitui decisão ética válida por si só. \citep{SU08}.


A orientação metodológica organiza pôsteres em mensagem central, apresentação visual e comunicação oral ajustadas ao público, em vez de reproduzir todo o artigo no painel principal. Guia voltado à pesquisa em saúde, não experimento de validação de um template universal; usado como fundamento editorial anterior à janela de 2016–2026, sem inferir eficácia educacional. \citep{SU09}.
